\SourcePage{263}
\SGAsection{3}{群概形 $(P,Q)$ 由 $\mathbf G_m$ 给出的双扩张：一般性质}
本节固定一个基\Term{scheme}{概形} $S$，并按照第0节的约定，在 $S$ 的 fppf \Term{topos}{拓扑斯}中工作。

\textbf{3.1.} 首先回顾群 $\operatorname{Biext}^1(P,Q;\mathbf G_m)$ 最重要的例子：$P$ 和 $Q$ 都是 $S$ 上的\Term{abelian-scheme}{阿贝尔概形}，而 $\operatorname{Biext}^1$ \Term{solvable}{可解}释为 $(P,Q)$ 上的\Term{divisorial-correspondence}{除子对应}群（第VII讲2.9.4）。此时，上一节的 $\ell$ 进\Term{pairing}{配对}就是\Term{abelian-scheme}{阿贝尔概形}理论中熟知的\Term{pairing}{配对} [1]。

\textbf{3.2.} 更\Term{generic}{一般}地，只假设 $Q$ 是\Term{abelian-scheme}{阿贝尔概形}。众所周知（第VII讲1.3.8），有
\[
\underline{\operatorname{Hom}}(Q,\mathbf G_m)=0,\tag{3.2.1}
\]
因而由1.4得到\Term{canonical-isomorphism}{典范同构}
\[
\operatorname{Biext}^1(P,Q;\mathbf G_m)\xrightarrow{\sim}\operatorname{Hom}(P,Q^*),\tag{3.2.2}
\]
其中
\[
Q^*=\underline{\operatorname{Ext}}^1(Q,\mathbf G_m)\tag{3.2.3}
\]
是 $Q$ 的\Term{dual-abelian-scheme}{对偶阿贝尔概形}。（\Term{generic}{一般}情形下 $Q^*$ 的\Term{representability}{可表性}由 M. Raynaud 证明，参见 [1]。）
若 $P$ 也是\Term{abelian-scheme}{阿贝尔概形}，则与（3.2.2）\Term{symmetric}{对称}地，有\Term{canonical-isomorphism}{典范同构}
\[
\operatorname{Biext}^1(P,Q;\mathbf G_m)\xrightarrow{\sim}\operatorname{Hom}(Q,P^*).\tag{3.2.4}
\]
\SourcePage{264}
因此，给出一个 $(P,Q)$ 由 $\mathbf G_m$ 给出的\Term{biextension}{双扩张} $\mathscr E$，既\Term{solvable}{可解}释为给出一个\Term{homomorphism}{同态}
\[u:P\longrightarrow Q^*,\]
也\Term{solvable}{可解}释为给出一个\Term{homomorphism}{同态}
\[v:Q\longrightarrow P^*.\]
此外，已知\Term{abelian-scheme}{阿贝尔概形}的\Term{homomorphism}{同态} $w:A\to B$ 是\Term{isomorphism}{同构}，当且仅当对每个整数 $n>0$，诱导的\Term{homomorphism}{同态} ${}_nw:{}_nA\to{}_nB$ 是\Term{isomorphism}{同构}。对上述 $u,v$，${}_nu:{}_nP\to{}_nQ^*$ 等同于由\Term{pairing}{配对} $\varphi_{\mathscr E}^n:{}_nP\times{}_nQ\to\mathbf G_m$ 定义的\Term{homomorphism}{同态} ${}_nP\to D({}_nQ)$（参见2.3.11），而 ${}_nv:{}_nQ\to{}_nP^*$ 等同于由\Term{pairing}{配对} $\psi_{\mathscr E}^n:{}_nP\times{}_nQ\to\mathbf G_m$ 定义的\Term{homomorphism}{同态} ${}_nQ\to D({}_nP)$。由 $\varphi_{\mathscr E}^n=-\psi_{\mathscr E}^n$（2.1.10），${}_nu$ 与 ${}_nv$ 相差一个\Term{signature}{符号}后互为\Term{transpose}{转置}。又因 ${}_nP$ 和 ${}_nQ$ 在 $S$ 上\Term{finite-locally-free}{有限局部自由}，从而满足\Term{cartier-duality}{Cartier 对偶性}，可知 ${}_nu$ 是\Term{isomorphism}{同构}当且仅当 ${}_nv$ 是\Term{isomorphism}{同构}，因此 $u$ 是\Term{isomorphism}{同构}当且仅当 $v$ 是\Term{isomorphism}{同构}。

这一结果也可\Term{equivalence}{等价}地表述如下，虽然这种表述稍欠\Term{symmetric}{对称}：考虑 $(P,P^*)$ 由 $\mathbf G_m$ 给出的\Term{canonical}{典范}\Term{biextension}{双扩张}（\Term{weil-biextension}{Weil 双扩张}，对应于恒等\Term{homomorphism}{同态} $v:P^*\to P^*$），则相应的\Term{canonical}{典范}\Term{homomorphism}{同态} $u:P\to P^{**}$ 是\Term{isomorphism}{同构}。这就是 Barsotti 与 Cartier 的\Term{biduality-theorem}{双对偶定理}。

\SourcePage{265}
\textbf{3.3.} 现在转向 $Q$ 在 $S$ 上\Term{finite-locally-free}{有限局部自由}，或是 $S$ 上\Term{multiplicative-type}{乘法型}\Term{scheme}{概形}的情形。其关键\Term{lemma}{引理}是下面这个熟知的结果。
\begin{SGAresult}{命题}{3.3.1}
设 $Q$ 是 $S$ 上的\Term{commutative-group-scheme}{交换群概形}，并假设 $Q$ 在 $S$ 上\Term{finite-locally-free}{有限局部自由}，或者在 $S$ 上是\Term{finite-type}{有限型}且为\Term{multiplicative-type}{乘法型}（SGA 3，第IX讲）。则
\[\underline{\operatorname{Ext}}^1(Q,\mathbf G_m)=0.\]
\end{SGAresult}
\begin{proof}
分两种情形讨论。

 a）$Q$ 在 $S$ 上\Term{finite-locally-free}{有限局部自由}。须证明 $Q$ 由 $\mathbf G_m$ 给出的每个\Term{extension}{扩张} $E$ 都对 fppf \Term{topology}{拓扑}\Term{local-scheme}{局部}\Term{split-extension}{分裂}。显然可假设 $S$ 是\Term{affine}{仿射}的；于是存在整数 $n>0$ 使 $nQ=0$（SGA 3，第IX讲6.1）。另一方面，$n\cdot\mathrm{id}$ 对 fppf 是 $\mathbf G_m$ 中的\Term{epimorphism}{满态射}，故 Kummer \Term{exact-sequence}{正合列}
\[
0\longrightarrow\boldsymbol\mu_n\longrightarrow\mathbf G_m\xrightarrow{n}\mathbf G_m\longrightarrow0
\]
产生一个 $\underline{\operatorname{Ext}}^i(Q,-)$ 的\Term{exact-sequence}{正合列}，从而给出\Term{epimorphism}{满态射}
\[
\underline{\operatorname{Ext}}^1(Q,\boldsymbol\mu_n)\longrightarrow{}_n\underline{\operatorname{Ext}}^1(Q,\mathbf G_m)=\underline{\operatorname{Ext}}^1(Q,\mathbf G_m).
\]
最后一个等式来自这样一个事实：乘以 $n$ 在 $Q$ 中为零，因此在 $\underline{\operatorname{Ext}}^1(Q,\mathbf G_m)$ 中也为零。这表明\Term{extension}{扩张} $E$ 来自 $Q$ 由 $\boldsymbol\mu_n$ 给出的\Term{extension}{扩张} $F$；
\SourcePage{266}
由 fpqc \Term{descent}{下降}理论（SGA 1，第VII讲），$F$ 可由 $S$ 上的\Term{finite-locally-free}{有限局部自由}\Term{scheme}{概形}表示。根据\Term{finite-locally-free}{有限局部自由}群的\Term{cartier-duality}{Cartier 对偶性}的已知性质（SGA 3，第VII讲3），其\Term{dual}{对偶}
\[D(F)=\underline{\operatorname{Hom}}(F,\mathbf G_m)\]
是 $D(\boldsymbol\mu_n)$ 由 $D(Q)$ 给出的\Term{extension}{扩张}。（须说明，若 $A\to B$ 是 $S$ 上\Term{finite-locally-free}{有限局部自由}\Term{commutative-group}{交换群}的\Term{monomorphism}{单态射}，则 $D(B)\to D(A)$ 是\Term{faithfully-flat}{忠实平坦}\Term{morphism}{态射}。由于 $D(B)$ 和 $D(A)$ 在 $S$ 上也\Term{finite-locally-free}{有限局部自由}，可应用按\Term{fiber}{纤维}判定\Term{faithfully-flat}{忠实平坦}性的准则，将问题化为 $S$ 是域 $k$ 的\Term{spectrum}{谱}的情形。此时，断言就是\Term{affine}{仿射}环上的\Term{homomorphism}{同态} $k(D(A))\to k(D(B))$ 为\Term{injectivity}{单射}（SGA 3，第VI$_B$讲11.14）。但此\Term{homomorphism}{同态}是 $k(B)\to k(A)$ 的\Term{transpose}{转置}；因 $A\to B$ 是\Term{closed-immersion}{闭浸入}（SGA 3，第VI$_B$讲1.4.2），后者\Term{surjectivity}{满射}。因此 $k(D(A))\to k(D(B))$ 确为\Term{injectivity}{单射}，证毕。）

\Term{inclusion}{包含映射} $\boldsymbol\mu_n\hookrightarrow\mathbf G_m$ 等同于 $D(\boldsymbol\mu_n)=\underline{\operatorname{Hom}}(\boldsymbol\mu_n,\mathbf G_m)$ 的一个\Term{section}{截面}。于是只须证明，此\Term{inclusion}{包含映射}在 $S$ 上 fppf \Term{local-scheme}{局部}可延拓为\Term{homomorphism}{同态} $F\to\mathbf G_m$，也就是说，所考虑的 $D(\boldsymbol\mu_n)$ 的\Term{section}{截面}可 fppf \Term{local-scheme}{局部}\Term{lifting}{提升}为 $D(F)$ 的\Term{section}{截面}。由于 $D(F)\to D(\boldsymbol\mu_n)$ \Term{faithfully-flat}{忠实平坦}且为\Term{finite-presentation}{有限表示}，这点显然成立。

 b）根据 SGA 3，第X讲4.5，在 $S$ 上对\Term{etale-topology}{平展拓扑}\Term{localization}{局部化}后，可假设 $Q$ \Term{diagonalizable}{可对角化}，即形如 $D(M)$，其中 $M$ 是\Term{finite-type}{有限型} $\mathbf Z$ 模。设 $T$ 是 $M$ 的\Term{torsion-submodule}{挠子模}，$L=M/T$ 是\Term{quotient-module}{商模}，它\Term{isomorphism}{同构}于 $\mathbf Z^r$。
\SourcePage{267}
于是\Term{exact-sequence}{正合列} $0\to T\to M\to L\to0$ 给出\Term{exact-sequence}{正合列}
\[0\longrightarrow D(L)\longrightarrow D(M)\longrightarrow D(L)\longrightarrow0\]
（SGA 3，第IX讲3.1），故只须分别证明 $D(L)$、$D(T)$ 与 $\mathbf G_m$ 的 $\underline{\operatorname{Ext}}^1$ 为零。$D(T)$ 的情形可用 a）；$D(L)$ 的情形可化为 $L=\mathbf Z$，因此只须验证
\[\underline{\operatorname{Ext}}^1(\mathbf G_m,\mathbf G_m)=0,\]
即 $\mathbf G_m$ 由 $\mathbf G_m$ 给出的每个\Term{extension}{扩张} $E$ 都 fppf \Term{locally-trivial}{局部平凡}。这样的\Term{extension}{扩张}可由 $S$ 上\Term{affine}{仿射}且\Term{flat}{平坦}的\Term{group-scheme}{群概形}表示（fpqc \Term{descent}{下降}理论，SGA 1，第VII讲），其\Term{fiber}{纤维}都是\Term{torus}{环面}，所以它本身是\Term{torus}{环面}（SGA 3，第X讲4.9及\SGAsourcegap{原PDF第267页此处“et”之后的第二项引注留白，现有源未提供该引注。}）。根据 SGA 3，第X讲4.5，可假设 $E$ 是平凡\Term{torus}{环面} $D(R)\simeq\mathbf G_m^2$。再作 Zariski \Term{localization}{局部化}，可假设\Term{monomorphism}{单态射} $\mathbf G_m\hookrightarrow E=D(R)$ 来自一个\Term{constant}{常值}\Term{morphism}{态射} $R\to\mathbf Z_S$（SGA 3，第VIII讲1.4），而它必为\Term{epimorphism}{满态射}（SGA 3，第VIII讲3.2 b））。选择一个在 $\mathbf Z$ 的元素 $1$ 上方的 $r\in R$，便得到\Term{homomorphism}{同态} $E\to\mathbf G_m$，使\Term{extension}{扩张} $E$ \Term{split-extension}{分裂}，这就证明了所需断言。
\end{proof}
\begin{SGAresult}{推论}{3.3.2}
设 $Q$ 如3.3.1。则对每个 $P$，都有\Term{canonical-isomorphism}{典范同构}
\[
\operatorname{Biext}^1(P,Q;\mathbf G_m)\xleftarrow{\sim}\operatorname{Ext}^1(P,D(Q)),\tag{3.3.2.1}
\]
其中 $D(Q)$ 表示 $Q$ 的\Term{cartier-dual}{Cartier 对偶} $\underline{\operatorname{Hom}}(Q,\mathbf G_m)$。更确切地，\Term{canonical}{典范}\Term{functor}{函子}（1.1.6）给出 $P$ 由 $D(Q)$ 给出的\Term{extension}{扩张}\Term{category}{范畴}与 $(P,Q)$ 由 $\mathbf G_m$ 给出的\Term{biextension}{双扩张}\Term{category}{范畴}的\Term{equivalence-of-categories}{等价}。
\end{SGAresult}
\SourcePage{268}
\begin{proof}
结合3.3.1，这正是1.5的特例。
\end{proof}
\textbf{注记 3.3.3。} 在3.3.1的条件下，$D(Q)$ 可由 $S$ 上\Term{separated}{分离}、\Term{locally-finitely-presented}{局部有限表示}且\Term{locally-quasi-finite}{局部拟有限}的\Term{group-scheme}{群概形}表示。更确切地，若 $Q$ 在 $S$ 上\Term{finite-locally-free}{有限局部自由}，则 $D(Q)$ 亦然，这是熟知的（SGA 3，第VII讲3）；若 $Q$ 为\Term{multiplicative-type}{乘法型}，则 $D(Q)$ 是“\Term{finitely-generated-twisted-constant-group}{有限生成的扭常值群}”（SGA 3，第X讲5.1），这已在 SGA 3，第IX讲4.5及第X讲5.7中说明。因此，若 $P$ 是\Term{group-scheme}{群概形}，则根据 fpqc \Term{descent}{下降}理论中的一个熟知\Term{lemma}{引理}（SGA 3，第X讲5.4），$P$ 由 $D(Q)$ 给出的每个\Term{extension}{扩张}都是\Term{representable}{可表}的。

\textbf{3.3.4.} “对每个 $S$ 上\Term{affine}{仿射}、\Term{flat}{平坦}且\Term{finite-presentation}{有限表示}的\Term{group-scheme}{群概形} $Q$，都有 $\underline{\operatorname{Ext}}^1(Q,\mathbf G_m)=0$”这一断言是错误的，甚至对 $Q=\mathbf G_{a_S}$、$S=\operatorname{Spec}k[t]/(t^2)$ 也不成立，其中 $k$ 是\Term{characteristic}{特征} $p>0$ 的域。

\begin{SGAresult}{命题}{3.4}
设 $P$ 是 $S$ 上\Term{smooth}{光滑}、\Term{finite-presentation}{有限表示}且\Term{fiber}{纤维}\Term{connected}{连通}的\Term{group-scheme}{群概形}，$Q$ 是 $S$ 上的\Term{torus}{环面}。则 $(P,Q)$ 由 $\mathbf G_m$ 给出的\Term{biextension}{双扩张}\Term{category}{范畴}\Term{equivalence-of-categories}{等价}于\Term{point-category}{单点范畴}，即
\[
\operatorname{Biext}^0(P,Q;\mathbf G_m)=\operatorname{Biext}^1(P,Q;\mathbf G_m)=0.\tag{3.4.1}
\]
\end{SGAresult}
\begin{proof}
令 $M=D(Q)$，它是 $S$ 上\Term{fiber}{纤维}\Term{isomorphism}{同构}于 $\mathbf Z^r$ 的\Term{twisted-constant-group}{扭常值群}（3.3.3）。由3.3.2，所考虑的\Term{category}{范畴}\Term{equivalence-of-categories}{等价}于 $P$ 由 $M$ 给出的\Term{extension}{扩张}\Term{category}{范畴}。这个\Term{category}{范畴}是\Term{rigid-category}{刚性的}，即每个\Term{object}{对象}的\Term{automorphism-group}{自同构群}
\SourcePage{269}
都只有\Term{identity-element}{单位元}，这正是（3.4.1）的第一个等式，也可写成 $\operatorname{Hom}(P,M)=0$。而这个等式直接来自 $M$ 在 $S$ 上\Term{etale}{平展}且 $P$ 的\Term{fiber}{纤维}\Term{connected}{连通}。所以只须证明
\[\operatorname{Ext}^1(P,M)=0,\]
即 $P$ 由 $M$ 给出的每个\Term{extension}{扩张} $E$ 都平凡。利用此\Term{extension}{扩张}的\Term{splitting}{分裂截面}在存在时的唯一性，只须证明 $E$ \Term{locally-trivial}{局部平凡}，从而化为 $M=\mathbf Z_S^r$，甚至 $M=\mathbf Z_S$ 的情形。若 $S$ 是域的\Term{spectrum}{谱}，结论来自下文5.5（i）。（注意：我们在第IX讲6之前不会使用3.4，因而尤其没有循环论证！）现说明如何把\Term{generic}{一般}情形化为这个特例。只须找出 $P$ 上 $P$ 的一个\Term{section}{截面} $f$，使 $f(0)=0$，因为这样的\Term{section}{截面}必为\Term{multiplication}{乘法}性的，即满足
\[f(x+y)-f(x)-f(y)=0.\]
事实上，此等式左端表示一个\Term{morphism}{态射} $P\times_SP\to\mathbf Z_S$，它在 $P\times_SP$ 的零\Term{section}{截面}上为零，而 $P\times_SP$ 的\Term{fiber}{纤维}\Term{connected}{连通}（因为 $P$ 的\Term{geometric-fiber}{几何纤维}\Term{connected}{连通}），故处处为零。于是所需结论来自下面的\Term{lemma}{引理}。
\end{proof}
\begin{SGAresult}{\Term{lemma}{引理}}{3.4.2}
设 $S$ 是\Term{scheme}{概形}，$P$ 是 $S$ 上\Term{smooth}{光滑}、\Term{finite-presentation}{有限表示}且\Term{geometric-fiber}{几何纤维}\Term{connected}{连通}的\Term{scheme}{概形}，$Z$ 是群，$E$ 是一个 $Z_P$ \Term{torsor}{挠子}（对 fppf \Term{topology}{拓扑}、\Term{etale-topology}{平展拓扑}或 fpqc \Term{topology}{拓扑}而言都一样，参见 SGA 3，第X讲5.4等），$e$ 是 $P$ 在 $S$ 上的\Term{section}{截面}，$e'$ 是 $E$ 在 $S$ 上位于 $e$ 上方的\Term{section}{截面}（即 $e^{-1}(E)$ 在 $S$ 上的\Term{section}{截面}）。
\SourcePage{270}
则 $E$ 平凡的（必要且）充分条件是：对每个 $s\in S$，诱导的 $P_s$ 上的\Term{torsor}{挠子} $E_s$ 都平凡。此时存在唯一的 $E$ 在 $P$ 上的\Term{section}{截面} $f$，使 $f\circ e=e'$。
\end{SGAresult}
\begin{proof}
此\Term{section}{截面}的唯一性是清楚的：两个这样的\Term{section}{截面} $f,g$ 的“差”是一个\Term{morphism}{态射} $h:P\to Z_S$，满足 $h\circ e=1$；由于 $P$ 的\Term{fiber}{纤维}\Term{connected}{连通}，这个\Term{morphism}{态射}恒等于 $1$。此外，若 $E$ 在 $P$ 上存在\Term{section}{截面} $f$，总可用一个\Term{morphism}{态射} $P\to Z_S$（来自一个\Term{morphism}{态射} $z:S\to Z_S$）修正它，得到满足 $f'\circ e=e'$ 的\Term{section}{截面} $f'$；只须选择 $z$，使 $f(e(t))=e'(t)z(t)$。将这一观察用于每个 $s\in S$ 上方的\Term{fiber}{纤维}，便为每个点 $s\in S$ 得到唯一确定的 $E_s$ 在 $P_s$ 上的\Term{section}{截面} $f_s$，它把 $e_s$ 送到 $e'_s$。只须证明这些 $f_s$ 来自一个 $E$ 在 $P$ 上的\Term{section}{截面} $f$；\Term{equivalence}{等价}地（EGA IV，17.9.3），须证明各 $f_s(P_s)$ 的并集 $U\subset E$ 是 $E$ 的\Term{open-subset}{开子集}。

先假设 $S$ \Term{noetherian}{诺特}。利用如下事实，容易验证 $U$ \Term{constructible}{可构造}：对每个 $s\in S$，在赋予诱导\Term{reduced}{约化}结构的 $T=\overline{\{s\}}$ 中，存在 $s$ 的一个\Term{open-neighborhood}{开邻域} $T'$，以及 $E_{T'}$ 在 $P_{T'}$ 上的一个\Term{section}{截面}，它延拓 $f_s$，并与指定的\Term{section}{截面} $e,e'$ 相容（这直接来自 EGA IV，8.8.2）。根据 EGA IV，1.10.1的准则，尚须证明 $U$ 对\Term{generization}{泛化}\Term{stable}{稳定}。为此，用借助 EGA II，7.1.9的通常方法，可化为 $S$ 是\Term{discrete-valuation-ring}{离散赋值环}的\Term{spectrum}{谱}的情形。由于 $P$ 在 $S$ 上\Term{smooth}{光滑}，
\SourcePage{271}
$P$ \Term{regular}{正则}，从而\Term{normal}{正规}；又因其\Term{fiber}{纤维}\Term{connected}{连通}，$P$ \Term{connected}{连通}。于是 $E$ 在 $P$ 相对于 $S$ 的\Term{generic-fiber}{泛纤维}上给定的\Term{section}{截面}，唯一延拓为 $E$ 在 $P$ 上的\Term{section}{截面} $f$（EGA IV，18.10.19和18.10.20），显然与 $e,e'$ 相容。因此，此时 $U=f(P)$ 是开的。这就在 $S$ \Term{noetherian}{诺特}时证明了3.4。\Term{generic}{一般}情形容易利用通常的取\Term{limit}{极限}方法及 EGA IV，8、9化为这一情形，细节留给读者。
\end{proof}
\textbf{注记 3.4.3。} 也可用几条高深的文献引用来“压倒”\Term{lemma}{引理}3.4.2，即 SGA 4，第XV讲1.15的 $n=1$ 情形，以及4.1。

\begin{SGAresult}{推论}{3.5}
设 $P$ 是 $S$ 上\Term{smooth}{光滑}、\Term{finite-presentation}{有限表示}且\Term{fiber}{纤维}\Term{connected}{连通}的\Term{group-scheme}{群概形}，$0\to Q'\to Q\to Q''\to0$ 是 $S$ 上\Term{commutative-group}{交换群}的\Term{exact-sequence}{正合列}，其中 $Q'$ 是\Term{torus}{环面}。则“\Term{biextension}{双扩张}的\Term{inverse-image}{逆像}”\Term{functor}{函子}
\[
\operatorname{BIEXT}(P,Q'';\mathbf G_m)\xrightarrow{\sim}\operatorname{BIEXT}(P,Q;\mathbf G_m)\tag{3.5.1}
\]
是\Term{equivalence-of-categories}{范畴等价}；特别地，以下自然\Term{homomorphism}{同态}都是\Term{isomorphism}{同构}：
\[
\operatorname{Biext}^i(P,Q'';\mathbf G_m)\xrightarrow{\sim}\operatorname{Biext}^i(P,Q;\mathbf G_m),\qquad i=0,1.\tag{3.5.2}
\]
\end{SGAresult}
\begin{proof}
在\Term{equivalence-of-categories}{等价}意义下，$\operatorname{BIEXT}(P,Q'';\mathbf G_m)$ 等同于如下\Term{category}{范畴}：\Term{object}{对象}是 $(P,Q)$ 由 $\mathbf G_m$ 给出的\Term{biextension}{双扩张} $E$，并给诱导的 $(P,Q')$ 由 $\mathbf G_m$ 给出的\Term{biextension}{双扩张} $E'$ 配备一个\Term{trivialization}{平凡化}（第VII讲3.7.6）。由3.4，对给定的\Term{biextension}{双扩张}，$E'$ 恰有一个\Term{trivialization}{平凡化}。当然，也可
\SourcePage{272}
利用与\Term{exact-sequence}{正合列} $0\to Q'\to Q\to Q''\to0$ 相联系的 $\operatorname{Biext}^i$ \Term{exact-sequence}{正合列}（第VII讲3.7.5），证明 $i=0,1$ 时 $\operatorname{Biext}^i(P,Q'';\mathbf G_m)\to\operatorname{Biext}^i(P,Q;\mathbf G_m)$ 为\Term{isomorphism}{同构}，这与3.5的结论\Term{equivalence}{等价}。
\end{proof}
\textbf{注记 3.6。} 当然，将3.5中的左右交换所得的\Term{symmetric}{对称}命题也成立；通过取\Term{symmetric-biextension}{对称双扩张}的运算（第VII讲2.7），它在\Term{form}{形式}上也可由3.5\Term{pushout}{推出}。在3.5的数据和假设之外，再假设具有\Term{symmetric}{对称}的数据和假设，即 $Q$ 的\Term{fiber}{纤维}\Term{connected}{连通}（或 $Q''$ 的\Term{fiber}{纤维}\Term{connected}{连通}；因 $Q'$ 的\Term{fiber}{纤维}\Term{connected}{连通}，二者\Term{equivalence}{等价}），并有\Term{exact-sequence}{正合列} $0\to P'\to P\to P''\to0$，其中 $P'$ 是\Term{torus}{环面}。依次应用3.5，以及将其\Term{symmetric}{对称}命题用于这个\Term{exact-sequence}{正合列}和 $Q''$，可知\Term{functor}{函子}
\[
\operatorname{BIEXT}(P'',Q'';\mathbf G_m)\xrightarrow{\sim}\operatorname{BIEXT}(P,Q;\mathbf G_m)\tag{3.6.1}
\]
是\Term{equivalence-of-categories}{范畴等价}；特别地，以下\Term{homomorphism}{同态}为\Term{isomorphism}{同构}：
\[
\operatorname{Biext}^i(P'',Q'';\mathbf G_m)\xrightarrow{\sim}\operatorname{Biext}^i(P,Q;\mathbf G_m),\qquad i=0,1.\tag{3.6.2}
\]

\begin{SGAresult}{命题}{3.7}
设 $A$ 是 $S$ 上的\Term{abelian-scheme}{阿贝尔概形}，$T$ 是 $S$ 上\Term{finite-type}{有限型}的\Term{multiplicative-type}{乘法型}\Term{group-scheme}{群概形}（SGA 3，第IX讲），$M=\underline{\operatorname{Hom}}(T,\mathbf G_m)$ 是 $T$ 的\Term{cartier-dual}{Cartier 对偶}（它是 $S$ 上\Term{finitely-generated-twisted-constant-group}{有限生成的扭常值群}，见3.3.3）。则分别由 $A$ 由 $T$ 给出的\Term{commutative-extension}{交换扩张}和 $(A,M)$ 由 $\mathbf G_m$ 给出的\Term{biextension}{双扩张}组成的\Term{category}{范畴} $\operatorname{EXT}(A,T)$ 与 $\operatorname{BIEXT}(A,M;\mathbf G_m)$ 都是\Term{rigid-category}{刚性的}，并\Term{canonical}{典范}\Term{equivalence-of-categories}{等价}；
\SourcePage{273}
这些\Term{category}{范畴}的\Term{object}{对象}类组成的群，都\Term{canonical}{典范}\Term{isomorphism}{同构}于 $\operatorname{Hom}(M,A^*)$，其中 $A^*=\underline{\operatorname{Ext}}^1(A,\mathbf G_m)$ 是 $A$ 的\Term{dual-abelian-scheme}{对偶阿贝尔概形}（3.2）。特别地，有\Term{canonical-isomorphism}{典范同构}
\[
\operatorname{Ext}^1(A,T)\xrightarrow{\sim}\operatorname{Hom}(M,A^*).\tag{3.7.1}
\]
\end{SGAresult}
\begin{proof}
这正是在1.6中取 $P=A,Q=M$ 所得的特例。由\Term{biduality-theorem}{双对偶定理}（SGA 3，第X讲5.7），确有 $T\simeq\underline{\operatorname{Hom}}(M,\mathbf G_m)$，而只须验证条件（1.6.1），这里它们写为
\[
\underline{\operatorname{Hom}}(A,\mathbf G_m)=0,\qquad
\operatorname{Hom}\bigl(A,\underline{\operatorname{Ext}}^1(M,\mathbf G_m)\bigr)=0.
\]
第一个等式已在（3.2.1）中回顾；第二个可加强为
\[\underline{\operatorname{Ext}}^1(M,\mathbf G_m)=0.\]
为验证后式，因问题在 $S$ 上是\Term{local-scheme}{局部}的，可假设 $M$ 是\Term{constant}{常值}的，等于 $M_S$；将 $M$ \Term{resolution}{分解}为循环 $\mathbf Z$ 模之和后，化为 $M=\mathbf Z$ 的情形（这是显然的），以及 $M=\mathbf Z/n\mathbf Z$ 的情形，后者来自 $\mathbf G_m$ 为 $n$ \Term{divisible-group}{可除群}（或较为高深地，来自3.3.1）。

注意，当 $T=\mathbf G_m$，从而 $M=\mathbf Z_S$ 时，（3.7.1）化为恒等\Term{isomorphism}{同构}（3.2.3）。当然，（3.7.1）的直接证明几乎是显然的，无须引入群
\[\operatorname{Ext}^1\bigl(A\mathbin{\overset{\mathbf L}{\otimes}}M,\mathbf G_m\bigr)=\operatorname{Biext}^1(A,M;\mathbf G_m).\]
\end{proof}

\SourcePage{274}
\SGAsection{4}{域上光滑连通群概形的扩张与双扩张}
回顾下面这个归于 Rosenlicht 的\Term{lemma}{引理} [3 bis，命题3]。
\begin{SGAresult}{\Term{lemma}{引理}}{4.1}
设 $k$ 是域，$X,Y$ 是两个\Term{finite-type}{有限型}、\Term{separable}{可分}且\Term{geometrically-connected}{几何连通}的 $k$ \Term{scheme}{概形}，分别配备 $k$ \Term{rational-point}{有理点} $e_X,e_Y$。则每个在两个\Term{subscheme}{子概形} $X\times e_Y$ 和 $e_X\times Y$ 上平凡的\Term{morphism}{态射} $X\times Y\to\mathbf G_m$ 都是平凡的。
\end{SGAresult}
我们只在 $X,Y$ 为 $k$ 上的\Term{group-scheme}{群概形}时使用这个结果；此时由假设可知，$X,Y$ 甚至在 $k$ 上\Term{smooth}{光滑}。注意，利用归纳法，4.1蕴含任意多个因子的类似结论。因此可以应用第VII讲1.3及第VII讲2.9的结果，得到以下两个命题。

\begin{SGAresult}{命题}{4.2}
设 $P$ 是域 $k$ 上的\Term{smooth}{光滑}\Term{connected}{连通}群。则从 $P$ 由 $\mathbf G_m$ 给出的群\Term{extension}{扩张}\Term{category}{范畴}到 $P$ 上在 $P$ 的\Term{identity-element}{单位元} $e_P$ 处\Term{rigidified}{刚化}的\Term{invertible-module}{可逆模}\Term{category}{范畴}的自然\Term{functor}{函子}是\Term{fully-faithful}{全忠实}的。其\Term{essential-image}{本质像}由满足
\[\pi^*(P)\simeq\operatorname{pr}_1^*(P)\operatorname{pr}_2^*(P)\]
的\Term{rigidified}{刚化}\Term{invertible-module}{可逆模} $\mathscr L$ 组成，其中 $\pi:P\times P\to P$ 是 $P$ 的\Term{composition-law}{合成律}，$\operatorname{pr}_1,\operatorname{pr}_2$ 是乘积的两个\Term{projection}{投影}。
\end{SGAresult}
\begin{SGAresult}{命题}{4.3}
设 $P,Q$ 是域 $k$ 上的两个\Term{smooth}{光滑}\Term{connected}{连通}\Term{algebraic-group}{代数群}。则从 $(P,Q)$ 由 $\mathbf G_m$ 给出的\Term{biextension}{双扩张}\Term{category}{范畴}到 $P\times Q$ 上相对于 $P,Q$ 的\Term{unit}{单位}\Term{section}{截面} $e_P,e_Q$ \Term{birigidified}{双刚化}的\Term{invertible-module}{可逆模}\Term{category}{范畴}的自然\Term{functor}{函子}
\SourcePage{275}
是\Term{fully-faithful}{全忠实}的，而且这两个\Term{category}{范畴}都是\Term{rigid-category}{刚性的}。$P\times Q$ 上的\Term{birigidified}{双刚化}\Term{invertible-module}{可逆模} $\mathscr L$ 来自一个\Term{biextension}{双扩张}，当且仅当它沿 $\pi_P\times\mathrm{id}_Q$ 和 $\mathrm{id}_P\times\pi_Q$ 在 $P\times P\times Q$ 和 $P\times Q\times Q$ 上的\Term{inverse-image}{逆像}满足“\Term{theorem-of-the-cube}{立方体定理}”（[1]或[2]）。
\end{SGAresult}
\begin{proof}
最后这个条件，正是第VII讲2.9.3 b）中的条件在第VII讲3.5 b）所给具体描述下的翻译。
\end{proof}
\textbf{4.4.} 我不知道4.3中的\Term{cube-condition}{立方体条件}是否自动满足。至少已知 [2，IV 2.6]，对给定的 $\mathscr L$，总存在一个整数 $n\geq1$，它是 $k$ 的\Term{characteristic-exponent}{特征指数}的幂，使 $\mathscr L^{\otimes n}$ 满足所考虑的条件，因而来自 $(P,Q)$ 由 $\mathbf G_m$ 给出的\Term{biextension}{双扩张}；并且总能找到 $k$ 的一个\Term{radicial-extension}{根性扩张} $k'$，使所需条件在\Term{base-extension}{基扩张} $k\to k'$ 后成立，从而这次\Term{group-extension}{扩张}后 $\mathscr L$ 给出一个来自 $(P,Q)$ 由 $\mathbf G_m$ 给出的\Term{biextension}{双扩张}的\Term{birigidified}{双刚化}\Term{invertible-module}{可逆模}。特别地，若 $k$ \Term{perfect-field}{完美}，则4.3中的条件总成立。由迹的论证，最后这个事实在\Term{form}{形式}上蕴含前面的结论；也可应用下面的4.7得到它，参见4.9。

把 EGA Err$_{\mathrm{IV}}$ 53，21.4.13的下述结果作为\Term{lemma}{引理}回顾如下。
\SourcePage{276}
\begin{SGAresult}{\Term{lemma}{引理}}{4.5}
设 $f:X\to Y$ 是\Term{smooth}{光滑}\Term{surjective-morphism}{满}\Term{morphism}{态射}，其\Term{fiber}{纤维}\Term{integral}{整}，$Y$ \Term{normal}{正规}且\Term{integral}{整}，\Term{generic-point}{泛点}为 $y$；$\mathscr L$ 是 $X$ 上的\Term{invertible-module}{可逆模}，其在 $X_y$ 上的\Term{restriction}{限制}平凡。则存在 $Y$ 上的\Term{invertible-module}{可逆模} $M$，使 $\mathscr L$ \Term{isomorphism}{同构}于 $f^*(M)$。
\end{SGAresult}
\begin{SGAresult}{命题}{4.6}
设 $P,Q$ 是域 $k$ 上两个\Term{smooth}{光滑}\Term{connected}{连通}\Term{algebraic-group}{代数群}，且 $P$ \Term{affine}{仿射}。再假设 $P$ 有一条\Term{composition-series}{合成列}，每个因子或为\Term{torus}{环面}，或\Term{isomorphism}{同构}于\Term{additive-group}{加法群} $\mathbf G_a$（当 $k$ \Term{perfect-field}{完美}时，这个条件尤其成立，见 SGA 3，第XVII讲7.2.1、4.1.1（i）$\Longleftrightarrow$（ii）及4.1.5）；或者假设 $Q$ 是\Term{abelian-variety}{阿贝尔簇} $B$ 由满足前述条件的\Term{smooth}{光滑}\Term{connected}{连通}\Term{affine}{仿射}\Term{algebraic-group}{代数群}给出的\Term{extension}{扩张}。在这些条件下，$P\times Q$ 上的每个\Term{birigidified}{双刚化}\Term{invertible-module}{可逆模}都平凡；因此（4.3），$(P,Q)$ 由 $\mathbf G_m$ 给出的每个\Term{biextension}{双扩张}更是平凡的。
\end{SGAresult}
\begin{proof}
首先，由于\Term{birigidified}{双刚化}模的每个\Term{automorphism}{自同构}都是恒等（4.3），\Term{descent}{下降}理论表明，只须证明 $\mathscr L$ 在 $k$ 的某个\Term{finite}{有限}\Term{separable-extension}{可分扩张}后平凡；取\Term{limit}{极限}后，这立即将我们化为 $k$ \Term{separably-closed}{可分闭}的情形。此时 $k$ 上的每个\Term{torus}{环面}都\Term{split-torus}{分裂}，所以若 $P$ 有4.6中所述的\Term{composition-series}{合成列}，它也有一条商\Term{isomorphism}{同构}于 $\mathbf G_a$ 或 $\mathbf G_m$ 的\Term{composition-series}{合成列}。反复应用4.5，并利用对每个域 $K$ 均有 $\operatorname{Pic}(\mathbf G_{a_K})=\operatorname{Pic}(\mathbf G_{m_K})=0$，得到
\[\operatorname{Pic}(P)=0.\]
\SourcePage{277}
这个等式在\Term{base-field}{基域} $k$ 的任意变换后仍成立。因此，将4.5用于 $\operatorname{pr}_2:P\times Q\to Q$ 和模 $\mathscr L$，可知 $\mathscr L$ \Term{isomorphism}{同构}于一个形如 $\operatorname{pr}_2^*(M)$ 的模，其中 $M$ 是 $Q$ 上的\Term{invertible-module}{可逆模}。由于 $\mathscr L$ 在 $e_P\times Q$ 上的\Term{restriction}{限制}平凡，立即得到 $M$ 平凡，故 $\mathscr L$ 平凡。这就证明了第一种情形的结论。

在第二种情形中，$Q$ 是\Term{abelian-scheme}{阿贝尔概形} $B$ 由群 $M$ 给出的\Term{extension}{扩张}，且 $M$ 有一条商\Term{isomorphism}{同构}于 $\mathbf G_m$ 或 $\mathbf G_a$ 的\Term{composition-series}{合成列}。将4.5用于\Term{morphism}{态射}
\[f:X=P\times Q\longrightarrow Y=P\times B.\]
它在 $P\times B$ 的\Term{generic-point}{泛点} $y$ 处的\Term{fiber}{纤维}是 $M_y$ 下的一个\Term{torsor}{挠子}；由 $M$ 的特殊结构及对每个域 $K$ 都成立的经典等式 $\mathrm H^1(K,\mathbf G_a)=0$、$\mathrm H^1(K,\mathbf G_m)=0$，这个\Term{torsor}{挠子}必平凡（这里取 $K=k(y)$）。由于 $\operatorname{Pic}(M)=0$，从而 $\operatorname{Pic}(X_y)=0$，的确可应用4.5，得到 $\mathscr L$ \Term{isomorphism}{同构}于 $Y=P\times B$ 上某个\Term{invertible-module}{可逆模} $\mathscr M$ 的\Term{inverse-image}{逆像}。将 $\mathscr M$ 换为
\[\mathscr M\operatorname{pr}_1^*(\mathscr M_1)^{-1}\operatorname{pr}_2^*(\mathscr M_2)^{-1},\]
其中 $\operatorname{pr}_1,\operatorname{pr}_2$ 是 $Y=P\times A$ 的两个\Term{projection}{投影}，$\mathscr M_1,\mathscr M_2$ 分别是 $\mathscr M$ 在 $P\times e_B$、$e_P\times B$ 上的\Term{restriction}{限制}，便可假设 $\mathscr M$ \Term{birigidified}{双刚化}，而 $\mathscr L$ 是它的\Term{birigidified}{双刚化}\Term{inverse-image}{逆像}。这将我们化为 $Q$ 本身就是\Term{abelian-scheme}{阿贝尔概形}的情形。此时 $\mathscr L$ 定义了一个\Term{canonical}{典范}\Term{morphism}{态射}
\[P\longrightarrow\underline{\operatorname{Pic}}_{Q/k},\tag{*}\]
它把\Term{identity-element}{单位元}送到\Term{identity-element}{单位元}。因 $P$ \Term{connected}{连通}，此\Term{morphism}{态射}
\SourcePage{278}
经过\Term{dual-abelian-scheme}{对偶阿贝尔概形} $Q'=\underline{\operatorname{Pic}}^0_{Q/k}$ \Term{resolution}{分解}。众所周知，从\Term{smooth}{光滑}\Term{connected}{连通}\Term{affine}{仿射}\Term{algebraic-group}{代数群}到\Term{abelian-scheme}{阿贝尔概形}的\Term{scheme}{概形}\Term{morphism}{态射}是\Term{constant}{常值}的。（这是个“几何”问题，即可假设 $k$ \Term{algebraically-closed}{代数闭}，此时通过逐层\Term{resolution}{分解}，立即化为 $P=\mathbf G_a$ 或 $\mathbf G_m$；而已知，从 $\mathbf P^1_k$ 到 $k$ 上\Term{abelian-scheme}{阿贝尔概形}的每个\Term{rational-map}{有理映射}都是\Term{constant}{常值}的。）由（*）为\Term{constant}{常值}，因而为零，可知 $\mathscr L$ \Term{isomorphism}{同构}于 $P$ 上某个\Term{invertible-module}{可逆模} $M$ 的\Term{inverse-image}{逆像}；又因 $\mathscr L$ 在 $P\times e_A$ 上平凡，$M$ 平凡，故 $\mathscr L$ 也平凡，证毕。
\end{proof}
\begin{SGAresult}{命题}{4.7}
设 $P,Q$ 是域 $k$ 上两个\Term{smooth}{光滑}\Term{connected}{连通}\Term{algebraic-group}{代数群}。假设 $P$ 是\Term{abelian-scheme}{阿贝尔概形} $A$ 由一个\Term{smooth}{光滑}\Term{connected}{连通}\Term{affine}{仿射}\Term{algebraic-group}{代数群} $L$ 给出的\Term{extension}{扩张}，并且 $L$ 有一条\Term{composition-series}{合成列}，其每个因子都\Term{isomorphism}{同构}于\Term{torus}{环面}或 $\mathbf G_a$（当 $k$ \Term{perfect-field}{完美}时，这个条件总成立）。则1.3中所考虑的\Term{functor}{函子}是\Term{equivalence-of-categories}{范畴等价}，即 $P\times Q$ 上的每个\Term{birigidified}{双刚化}\Term{invertible-module}{可逆模} $\mathscr L$ 都来自 $(P,Q)$ 的一个\Term{biextension}{双扩张}。此外，\Term{rigid-category}{刚性}\Term{category}{范畴}的\Term{functor}{函子}
\[\operatorname{BIEXT}(A,Q;\mathbf G_m)\longrightarrow\operatorname{BIEXT}(P,Q;\mathbf G_m)\]
是\Term{equivalence-of-categories}{范畴等价}，换言之，它给出群的\Term{isomorphism}{同构}
\[
\operatorname{Biext}^1(A,Q;\mathbf G_m)\xrightarrow{\sim}\operatorname{Biext}^1(P,Q;\mathbf G_m).\tag{4.7.1}
\]
\end{SGAresult}
\SourcePage{279}
\begin{proof}
对第一项断言，仍利用\Term{descent}{下降}和4.3化为 $k$ \Term{separably-closed}{可分闭}的情形。仿照4.6的证明，由4.5及对 $L$ 的假设，可知 $\mathscr L$ \Term{isomorphism}{同构}于 $A\times Q$ 上某个\Term{birigidified}{双刚化}模 $\mathscr M$ 的\Term{inverse-image}{逆像}；所以为证明4.7的第一项断言，只须考虑 $A=P$。此时 $\mathscr L$ 定义一个带\Term{base-point}{基点}的\Term{scheme}{概形}的\Term{morphism}{态射}
\[f:Q\longrightarrow\underline{\operatorname{Pic}}_{P/k},\]
它必经过 $P$ 的\Term{dual-abelian-scheme}{对偶阿贝尔概形} $P'$ \Term{resolution}{分解}。设 $\mathscr L'$ 是 $P\times P'$ 上的\Term{weil-sheaf}{Weil 层}沿
\[\mathrm{id}_P\times f:P\times Q\longrightarrow P\times P'\]
的\Term{inverse-image}{逆像}。它是\Term{birigidified}{双刚化}模；由 $f$ 的定义，$\mathscr L'$ 与 $\mathscr L$ 仅相差 $P$ 上某个\Term{invertible-module}{可逆模} $M$ 的\Term{inverse-image}{逆像}。由于 $\mathscr L'$ 和 $\mathscr L$ 在 $P\times e_Q$ 上的\Term{restriction}{限制}均平凡，$M$ 平凡，故 $\mathscr L$ \Term{isomorphism}{同构}于 $\mathscr L'$。我们已在第VII讲2.9.4中指出，两个\Term{abelian-scheme}{阿贝尔概形}的乘积上的\Term{divisorial-correspondence}{除子对应}都来自\Term{biextension}{双扩张}。因此 $\mathscr L'$，或\Term{equivalence}{等价}地 $\mathscr L$，来自一个\Term{biextension}{双扩张}。

回到\Term{generic}{一般}情形，即 $P$ 是 $A$ 由 $L$ 给出的\Term{extension}{扩张}，还须证明（4.7.1）确为\Term{isomorphism}{同构}。这来自与\Term{exact-sequence}{正合列}
\[0\longrightarrow L\longrightarrow P\longrightarrow A\longrightarrow0\]
相联系的 $\operatorname{Biext}^i$（$i=0,1$）的\Term{exact-sequence}{正合列}，以及等式
\[\operatorname{Biext}^0(L,Q;G)=0,\qquad\operatorname{Biext}^1(L,Q;G)=0.\]
\SourcePage{280}
第一个等式是4.3中\Term{rigid-category}{刚性}断言的特例，第二个正是4.6的第一种情形。
\end{proof}
\begin{SGAresult}{推论}{4.8}
在4.7的条件下，假设 $Q$ 是一个\Term{algebraic-group}{代数群} $B$（它必\Term{smooth}{光滑}且\Term{connected}{连通}）由一个\Term{smooth}{光滑}\Term{connected}{连通}\Term{affine}{仿射}\Term{algebraic-group}{代数群} $M$ 给出的\Term{extension}{扩张}。则\Term{rigid-category}{刚性}\Term{category}{范畴}的\Term{functor}{函子}
\[\operatorname{BIEXT}(A,B;\mathbf G_m)\longrightarrow\operatorname{BIEXT}(P,Q;\mathbf G_m)\]
是\Term{equivalence-of-categories}{范畴等价}，换言之，它给出群的\Term{isomorphism}{同构}
\[
\operatorname{Biext}^1(A,B;\mathbf G_m)\xrightarrow{\sim}\operatorname{Biext}^1(P,Q;\mathbf G_m).\tag{4.8.1}
\]
\end{SGAresult}
\begin{proof}
考虑（4.7.1）后，只须证明有\Term{isomorphism}{同构}
\[\operatorname{Biext}^1(A,B;\mathbf G_m)\xrightarrow{\sim}\operatorname{Biext}^1(A,Q;\mathbf G_m).\]
这与（4.7.1）一样，来自\Term{exact-sequence}{正合列}
\[0\longrightarrow M\longrightarrow Q\longrightarrow B\longrightarrow0\]
所诱导的 $\operatorname{Biext}^i$ 的\Term{exact-sequence}{正合列}以及等式
\[\operatorname{Biext}^0(P,M;\mathbf G_m)=0,\qquad\operatorname{Biext}^1(P,M;\mathbf G_m)=0,\]
其中第二个等式是4.6所讨论的第二种情形的特例（交换两个因子的\Term{action}{作用}）。
\end{proof}

\SourcePage{281}
\textbf{注记 4.9。} 设 $P,Q$ 是\Term{perfect-field}{完美}域 $k$ 上两个\Term{smooth}{光滑}\Term{connected}{连通}\Term{algebraic-group}{代数群}。已知 [3]，二者各为\Term{abelian-scheme}{阿贝尔概形}由\Term{smooth}{光滑}\Term{connected}{连通}\Term{affine}{仿射}群给出的\Term{extension}{扩张}，而后者是\Term{torus}{环面}与 $\mathbf G_a$ 群的多重\Term{extension}{扩张}之积（参见4.6所给引注）：
\[
0\longrightarrow L\longrightarrow P\longrightarrow A\longrightarrow0,\qquad
0\longrightarrow M\longrightarrow Q\longrightarrow B\longrightarrow0.
\]
所以可应用4.8；它表明，$(P,Q)$ 由 $\mathbf G_m$ 给出的\Term{biextension}{双扩张}理论，归结为\Term{abelian-variety}{阿贝尔簇}对 $(A,B)$ 的相应理论——这是多么经典的理论！

\begin{SGAresult}{推论}{4.10}
设 $P,Q$ 是 $k$ 上两个\Term{smooth}{光滑}\Term{connected}{连通}\Term{algebraic-group}{代数群}，$L$ 是 $P$ 的一个\Term{affine}{仿射}\Term{smooth}{光滑}\Term{connected}{连通}\Term{algebraic-subgroup}{代数子群}，$E$ 是 $(P,Q)$ 由 $\mathbf G_m$ 给出的\Term{biextension}{双扩张}，从而由2.2有从 $E$ \Term{derive}{导出}的\Term{pairing}{配对}
\[T_\ell(P)\times T_\ell(Q)\longrightarrow T_\ell(\mathbf G_m).\]
则诱导的\Term{pairing}{配对}
\[T_\ell(L)\times T_\ell(Q)\longrightarrow T_\ell(\mathbf G_m)\]
为零。
\end{SGAresult}
\begin{proof}
显然可假设 $k$ \Term{algebraically-closed}{代数闭}，此时满足4.6的条件，因而 $E$ 在 $(L,Q)$ 上的\Term{restriction}{限制}是平凡\Term{biextension}{双扩张}。所考虑的诱导\Term{pairing}{配对}就是这个\Term{restriction}{限制}所定义的\Term{pairing}{配对}，故结论成立。
\end{proof}

\SourcePage{282}
\textbf{4.11.} 设 $P,Q$ 是 $k$ 上两个\Term{smooth}{光滑}\Term{algebraic-group}{代数群}，$E$ 是 $(P,Q)$ 由 $\mathbf G_m$ 给出的\Term{biextension}{双扩张}，$T$ 是 $P$ 的最大子\Term{torus}{环面}。我们要确定\Term{pairing}{配对}
\[
T_\ell(P)\times T_\ell(Q)\longrightarrow T_\ell(\mathbf G_{m_k})\tag{4.11.1}
\]
的\Term{left-annihilator}{左零化子}，这里 $\ell$ 仍是一个固定\Term{prime-number}{素数}。由4.10，这个\Term{left-annihilator}{左零化子}无论如何都包含 $T_\ell(P)$ 的\Term{subgroup}{子群} $T_\ell(T)$。考虑 $k$ 的一个\Term{group-extension}{扩张} $k'$，使 $P,Q$ 在其上都可写成\Term{abelian-scheme}{阿贝尔概形}由\Term{affine}{仿射}\Term{group-scheme}{群概形}给出的\Term{extension}{扩张}，且后者是\Term{torus}{环面}和 $\mathbf G_a$ 群的多重\Term{extension}{扩张}。若 $P,Q$ 的\Term{unipotent-rank}{幺幂秩}为零，可取 $k'=k$；总之，可取 $k'$ 为 $k$ 的\Term{perfect-closure}{完美闭包}（或一个\Term{finite}{有限}\Term{radicial-extension}{根性扩张}）。设 $A'$、$B'$ 分别是 $P_{k'}$、$Q_{k'}$ 的\Term{abelian-part}{阿贝尔部分}。由4.7，$(P_{k'},Q_{k'})$ 的\Term{biextension}{双扩张} $E_{k'}$ 来自 $(A',B')$ 的\Term{biextension}{双扩张} $F'$，它在唯一\Term{isomorphism}{同构}意义下确定。因此，由（4.11.1）经\Term{base-field}{基域}变换 $k\to k'$ 得到的\Term{pairing}{配对}
\[
T_\ell(P_{k'})\times T_\ell(Q_{k'})\longrightarrow T_\ell(\mathbf G_{m_{k'}}),\tag{4.11.2}
\]
也就是 $(P_{k'},Q_{k'})$ 的\Term{biextension}{双扩张} $E_{k'}$ 所定义的\Term{pairing}{配对}，可由\Term{biextension}{双扩张} $F'$ 所定义的\Term{pairing}{配对}
\[
T_\ell(A')\times T_\ell(B')\longrightarrow T_\ell(\mathbf G_{m_{k'}})\tag{4.11.3}
\]
与\Term{canonical}{典范}\Term{homomorphism}{同态}
\[
T_\ell(P_{k'})\longrightarrow T_\ell(A'),\qquad T_\ell(Q_{k'})\longrightarrow T_\ell(B')\tag{4.11.4}
\]
复合而得（若 $P,Q$ 为 $\ell$ \Term{divisible-group}{可除群}，这些\Term{homomorphism}{同态}事实上是\Term{epimorphism}{满态射}）。

\SourcePage{283}
因此，\Term{pairing}{配对}（4.11.2）的\Term{left-annihilator}{左零化子}（它由（4.11.1）的\Term{left-annihilator}{左零化子}经\Term{base-change}{基变换}得到），正是\Term{pairing}{配对}（4.11.3）的\Term{left-annihilator}{左零化子}在（4.11.4）第一个\Term{homomorphism}{同态}下的\Term{inverse-image}{逆像}。若 $C'$ 是 $A'$ 的最大\Term{abelian-subscheme}{阿贝尔子概形}，使 $F'$ 在 $(C',B')$ 上的\Term{restriction}{限制}为平凡\Term{biextension}{双扩张}（因此 $C'$ 也是 $F'$ 定义的\Term{homomorphism}{同态}
\[
A'\longrightarrow B'\text{ 的对偶阿贝尔概形}\tag{4.11.5}
\]
的\Term{reduced}{约化}\Term{neutral-component}{单位分支}），则（4.11.3）的\Term{left-annihilator}{左零化子}等同于 $T_\ell(C')$；因而\Term{pairing}{配对}（1.11.1）的\Term{left-annihilator}{左零化子}，正是 $T_\ell(C')\hookrightarrow T_\ell(A')$ 在（4.11.4）第一个\Term{morphism}{态射}下的\Term{inverse-image}{逆像}。

特别地，说\Term{pairing}{配对}（4.11.1）的\Term{left-annihilator}{左零化子}仅为 $T_\ell(T)$，即\Term{pairing}{配对}（4.11.2）的\Term{left-annihilator}{左零化子}仅为（4.11.4）第一个\Term{morphism}{态射}的核 $T_\ell(T')$，就等于说\Term{pairing}{配对}（4.11.3）\Term{left-separating}{左分离}，即 $C'=0$；又等于说，$F'$ 定义的\Term{morphism}{态射}（4.11.5）的核\Term{finite}{有限}。最后这种表述说明，这个条件与 $\ell$ 无关；而无论用哪一种表述，显然它都与\Term{group-extension}{扩张} $k'$ 的选取无关。这个条件满足时，就称给定的 $(P,Q)$ 由 $\mathbf G_m$ 给出的\Term{biextension}{双扩张} $E$ 是\Term{left-separating-biextension}{左分离的}。由\Term{duality}{对偶性}可知，这也意味着 $F'$ 定义的\Term{morphism}{态射}
\[
B'\longrightarrow A'\text{ 的对偶阿贝尔概形}\tag{4.11.6}
\]
是\Term{epimorphism}{满态射}（即\Term{surjectivity}{满射}）。

\SourcePage{284}
\Term{symmetric}{对称}地定义\Term{right-separating-biextension}{右分离双扩张}：它对应于（4.11.6）的核\Term{finite}{有限}，即（4.11.5）为\Term{epimorphism}{满态射}的情形。若 $(P,Q)$ 由 $\mathbf G_{m_k}$ 给出的\Term{biextension}{双扩张} $E$ 既\Term{left-separating-biextension}{左分离}又\Term{right-separating-biextension}{右分离}，就称它\Term{separating-biextension}{分离}；这等于说（4.11.5）是\Term{isogeny}{同源}，也等于说（4.11.6）是\Term{isogeny}{同源}。

由定义显然可见，这些概念在\Term{base-field}{基域} $k$ 的任意\Term{group-extension}{扩张}下\Term{invariant}{不变}。

\textbf{4.12.} 设 $P$ 是\Term{smooth}{光滑}\Term{connected}{连通}\Term{group-scheme}{群概形}，$\mathscr L$ 是 $P$ 上的\Term{invertible-module}{可逆模}，令
\[
\delta(\mathscr L)=\pi^*(\mathscr L)\operatorname{pr}_1^*(\mathscr L)^{-1}\operatorname{pr}_2^*(\mathscr L)^{-1},\tag{4.12.1}
\]
记号如4.2。正如第VII讲1.3.4中已指出的，这样得到 $P\times P$ 上的一个\Term{birigidified}{双刚化}\Term{invertible-module}{可逆模}。由4.3，若 $\delta(\mathscr L)$ 来自一个\Term{biextension}{双扩张}，则后者在唯一\Term{isomorphism}{同构}意义下确定。特别地，由4.7，若 $P$ 是\Term{abelian-scheme}{阿贝尔概形} $A$ 由\Term{torus}{环面}和 $\mathbf G_a$ 群的多重\Term{extension}{扩张}给出的\Term{extension}{扩张}，则这个\Term{biextension}{双扩张}有定义。无论如何，只要有定义，它显然是\Term{symmetric-biextension}{对称双扩张}。在前述有利情形下，再假设 $P$ 的最大\Term{torus}{环面}\Term{split-torus}{分裂}，则可将 $\delta(\mathscr L)$ 的构造化为 $P$ 是\Term{abelian-scheme}{阿贝尔概形}的情形；只须注意，此构造显然对 $P$ 具有\Term{functoriality}{函子性}，而另一方面，由\Term{canonical}{典范}\Term{homomorphism}{同态}\Term{derive}{导出}的\Term{homomorphism}{同态}
\[
\operatorname{Pic}(A)\longrightarrow\operatorname{Pic}(P)\tag{4.12.2}
\]
\SourcePage{285}
——这里的\Term{canonical}{典范}\Term{homomorphism}{同态}为 $f:A\to P$——由4.5为\Term{surjectivity}{满射}（参见4.6证明开头）。因此可找到\Term{isomorphism}{同构}
\[\mathscr L^*\simeq f^*(\mathscr M),\]
其中 $\mathscr M$ 是 $\mathscr M$ 上的\Term{invertible-module}{可逆模}。于是 $\delta(\mathscr L)$“就是”$(A,A)$ 由 $\mathbf G_{m_k}$ 给出的\Term{biextension}{双扩张} $\delta(\mathscr M)$ 的\Term{inverse-image}{逆像}。

\begin{SGAresult}{命题}{4.13}
在4.12的条件下，假设 $\mathscr L$ 是 $P$ 上的\Term{ample-invertible-module}{丰沛可逆模}。若 $(P,P)$ 由 $\mathbf G_{m_k}$ 给出的\Term{biextension}{双扩张} $\delta(\mathscr L)$ 有定义，则它是\Term{separating-biextension}{分离的}（4.11）；若 $P$ 是\Term{abelian-scheme}{阿贝尔概形} $A$ 由满足4.7条件的\Term{affine}{仿射}\Term{smooth}{光滑}\Term{connected}{连通}群 $L$ 给出的\Term{extension}{扩张}，则 $\delta(\mathscr L)$ 所来自的 $(A,A)$ 的\Term{biextension}{双扩张}（由4.8）形如 $\delta(\mathscr M)$，其中 $\mathscr M$ 是 $A$ 上的\Term{ample-invertible-module}{丰沛可逆模}。
\end{SGAresult}
\begin{proof}
可化为 $k$ \Term{algebraically-closed}{代数闭}的情形。按定义，采用4.12的记号，此时只须证明 $(A,A)$ 的\Term{biextension}{双扩张} $\delta(\mathscr M)$ 是\Term{separating-biextension}{分离的}。M. Raynaud 的学位论文 [2，XI 1.11] 证明：若 $\mathscr L=f^*(\mathscr M)$ \Term{ample}{丰沛}，则 $\mathscr M$ 也\Term{ample}{丰沛}。由 A. Weil 的理论，已知这蕴含 $\delta(\mathscr M)$ \Term{separating-biextension}{分离} [1][2]，即从 $A$ 到其\Term{dual-abelian-scheme}{对偶阿贝尔概形}的\Term{homomorphism}{同态}是\Term{isogeny}{同源}。
\end{proof}

\textbf{注记 4.14。} 考虑 $P$ 是\Term{abelian-scheme}{阿贝尔概形} $A$ 由\Term{torus}{环面}和 $\mathbf G_a$ 群的多重\Term{extension}{扩张} $L$ 给出的\Term{extension}{扩张}这一有利情形，并考虑\Term{canonical}{典范}的复合\Term{homomorphism}{同态}
\SourcePage{286}
\[
\operatorname{NS}(A)\xrightarrow{\delta}\operatorname{Biext}^1(A,A;\mathbf G_m)\xrightarrow{\sim}\operatorname{Biext}^1(P,P;\mathbf G_m).\tag{4.14.1}
\]
由4.8，第二个\Term{arrow}{箭头}是\Term{bijection}{双射}。另一方面，在\Term{abelian-variety}{阿贝尔簇}理论中，已知第一个\Term{arrow}{箭头}诱导 $A$ 的\Term{neron-severi-group}{Néron--Severi 群} $\operatorname{NS}(A)$ 与 $\operatorname{Biext}^1(A,A;\mathbf G_m)=\operatorname{Corr}(A,A)$ 中由\Term{symmetric-biextension}{对称双扩张}（即\Term{divisorial-correspondence}{除子对应}）的类组成的\Term{subgroup}{子群}之间的\Term{isomorphism}{同构}。因此，复合（4.14.1）诱导 $\operatorname{NS}(A)$ 与 $(P,P)$ 由 $\mathbf G_m$ 给出的\Term{symmetric-biextension}{对称双扩张}群之间的\Term{isomorphism}{同构}（由4.7，这个群也可描述为 $(P,P)$ 上\Term{symmetric}{对称}的\Term{divisorial-correspondence}{除子对应}的类组成的群）。若 $(P,P)$ 上的一个\Term{divisorial-correspondence}{除子对应}来自群 $\operatorname{NS}(A)$ 的一个\Term{ample}{丰沛}元素，就应称之为\Term{ample-divisorial-correspondence}{丰沛除子对应}，并将相应的 $(P,P)$ 由 $\mathbf G_m$ 给出的\Term{biextension}{双扩张}称为\Term{ample-biextension}{丰沛双扩张}。这个条件显然在\Term{base-field}{基域}变换下\Term{invariant}{不变}，因此，只假设 $P$ \Term{smooth}{光滑}且\Term{connected}{连通}而不对 $P$ 作上述\Term{restriction}{限制}，前一概念仍然有意义。

再假设 $L$“在 $k$ 上\Term{solvable}{可解}”，即它可逐层\Term{resolution}{分解}为 $\mathbf G_a$ 和 $\mathbf G_m$ 群。由上述观察及映射（4.12.2）为\Term{surjectivity}{满射}，可知
\[
\operatorname{Pic}(P)\xrightarrow{\delta}\operatorname{Biext}^1(P,P;\mathbf G_m)\tag{4.14.2}
\]
的像由\Term{symmetric-biextension}{对称双扩张}组成，其核为 $\operatorname{Pic}(P)$ 的\Term{subgroup}{子群} $\operatorname{Pic}^0(P)$，即 $\operatorname{Pic}^0(A)$ 的像。因此，令
\[\operatorname{NS}(P)=\operatorname{Pic}(P)/\operatorname{Pic}^0(P)\qquad(\simeq\operatorname{NS}(A)),\]
\SourcePage{287}
则复合（4.14.1）也\Term{solvable}{可解}释为一个\Term{canonical-isomorphism}{典范同构}
\[
\operatorname{NS}(P)\xrightarrow{\sim}\operatorname{Biext}^1(P,P;\mathbf G_m)^{\mathrm{sym}}.\tag{4.14.3}
\]
最后，4.13中引用的 Raynaud \Term{theorem}{定理}可表述为：$\operatorname{Pic}(P)$ 的元素 $\mathscr E$ \Term{ample}{丰沛}，当且仅当它沿4.14.1的像 $\delta(\mathscr E)$ 是一个\Term{ample-biextension}{丰沛双扩张}类。

\SGAsection{5}{由无挠截断常值群概形给出的扩张与双扩张}
本节及下一节给出一些辅助结果，将在第7节用于 $Z=\mathbf Z$ 的情形。
\begin{SGAresult}{命题}{5.1}
设 $S$ 是\Term{geometrically-unibranch}{几何单支}且\Term{irreducible}{不可约}的\Term{scheme}{概形}，\Term{generic-point}{泛点}为 $\eta$，$Z$ 是\Term{torsion-free}{无挠}群，从而有 $S$ 上的\Term{constant-group-scheme}{常值群概形} $Z_S$。则\Term{group-scheme}{群概形} $Z_S$ 下的每个\Term{torsor}{挠子} $E$ 都平凡，而且从 $Z_S$ 下的\Term{torsor}{挠子}\Term{category}{范畴}到 $Z=Z_S(\eta)$ 下的\Term{torsor}{挠子}\Term{category}{范畴}的\Term{functor}{函子} $E\mapsto E(\eta)$ 是\Term{equivalence-of-categories}{范畴等价}。
\end{SGAresult}
\begin{proof}
\Term{descent}{下降}理论的一个熟知\Term{lemma}{引理}（SGA 3，第X讲5.4）表明，$G_S$ 下的每个\Term{torsor}{挠子} $E$ 都\Term{representable}{可表}（对 $S$ 上每个\Term{separated}{分离}、\Term{locally-finitely-presented}{局部有限表示}且\Term{locally-quasi-finite}{局部拟有限}的\Term{group-scheme}{群概形}，这个结果都成立）。由于 $Z_S$ 在 $S$ 上\Term{etale}{平展}，由 fpqc \Term{descent}{下降}，$Z_S$ 下的每个\Term{torsor}{挠子} $E$ 也在 $S$ 上\Term{etale}{平展}，因而对\Term{etale-topology}{平展拓扑}\Term{locally-trivial}{局部平凡}（不只是
\SourcePage{288}
对 fppf \Term{topology}{拓扑}\Term{locally-trivial}{局部平凡}）。所以，为证明 $E$ 平凡，只须证明
\[\mathrm H^1(S_{\text{ét}},Z_S)=e.\]
同样，5.1中\Term{fully-faithful}{全忠实}的断言化为等式
\[\mathrm H^0(S_{\text{ét}},Z_S)\xleftarrow{\sim}Z.\]
后一个等式显然成立，现证明前一个。考虑\Term{inclusion}{包含映射}
\[i:\eta\longrightarrow S.\]
由于 $S$ \Term{geometrically-unibranch}{几何单支}，有
\[i_*(Z_\eta)\xleftarrow{\sim}Z_S\]
（SGA 4，第IX讲2.14.1）。由替代 $i:\eta\to S$ 的 Leray \Term{spectral-sequence}{谱序列}的低次非交换\Term{exact-sequence}{正合列}（SGA 4，第XII讲3.2）
\[e\longrightarrow\mathrm H^1(S,i_*(Z_\eta))\longrightarrow\mathrm H^1(\eta,Z_\eta)\longrightarrow\cdots,\]
可知 $\mathrm H^1(S,Z_S)\to\mathrm H^1(\eta,Z_\eta)$ 是\Term{injectivity}{单射}。因此化为 $S$ 为域 $k$ 的\Term{spectrum}{谱}的情形。若 $\pi$ 是 $k$ 的一个\Term{separable-closure}{可分闭包}的\Term{galois-group}{Galois 群}，则
\[\mathrm H^1(k,Z)=\operatorname{Hom}_{\mathrm{cont}}(\pi,Z)/Z=e,\]
因为 $E$ \Term{torsion-free}{无挠}。证明完成。
\end{proof}
\begin{SGAresult}{推论}{5.2}
沿用5.1的记号，设 $P$ 是 $S$ 上的\Term{smooth}{光滑}\Term{scheme}{概形}。则 $Z_P$ 下的每个\Term{torsor}{挠子} $E$ 都平凡，而且从 $P$ 上以 $Z_P$ 为群的\Term{torsor}{挠子}\Term{category}{范畴}到
\SourcePage{289}
$P_\eta$ 上以 $Z_{P_\eta}$ 为群的\Term{torsor}{挠子}\Term{category}{范畴}的\Term{functor}{函子} $E\mapsto E_\eta$ 是\Term{equivalence-of-categories}{范畴等价}。
\end{SGAresult}
\begin{proof}
事实上，$P$ \Term{geometrically-unibranch}{几何单支}（EGA IV，17.5.7）；由于 $P$ 在 $S$ 上\Term{flat}{平坦}，$P$ 的每个\Term{irreducible-component}{不可约分支}都\Term{dominant}{支配} $S$，因而 $P$ 的\Term{irreducible-component}{不可约分支}族\Term{locally-finite}{局部有限}。所以 $P$ 是若干\Term{irreducible}{不可约}\Term{geometrically-unibranch}{几何单支}\Term{prescheme}{预概形} $P_i$ 的和。将5.1用于各个 $P_i$ 和 $P_{i\eta}$，并考虑到 $P_i$ 的\Term{generic-point}{泛点} $p_i$ 也是 $P_{i\eta}$ 的\Term{generic-point}{泛点}，便立即得到5.2。
\end{proof}
\textbf{5.2.1.} 上一推论立即给出 $P$ 上 $(Z_P,Z_P)$ 下的\Term{bitorsor}{双挠子}的类似结论：从这些\Term{bitorsor}{双挠子}的\Term{category}{范畴}到 $P_\eta$ 上 $(Z_{P_\eta},Z_{P_\eta})$ 下的\Term{bitorsor}{双挠子}\Term{category}{范畴}的\Term{functor}{函子} $E\mapsto E_\eta$ 是\Term{equivalence-of-categories}{范畴等价}。事实上，只须把 $E$ 的\Term{bitorsor}{双挠子}结构具体写成 $Z_P$ 下的右\Term{torsor}{挠子}结构，再配备一个从 $Z_P$ 到 $Z_P$ 在 $E$ 上\Term{action}{作用}的\Term{commutant}{交换子群}的\Term{isomorphism}{同构}；并利用 $E$ 必平凡，从而该\Term{commutant}{交换子群}\Term{isomorphism}{同构}于 $Z_P$ 本身。于是对 $P$ 的\Term{irreducible-component}{不可约分支} $X$，问题化为下面这个显然的观察：若 $X$ 是\Term{nonempty}{非空}\Term{connected}{连通}\Term{scheme}{概形}，则从\Term{set}{集合}\Term{category}{范畴}到 $X$ \Term{scheme}{概形}\Term{category}{范畴}的\Term{functor}{函子} $Z\mapsto Z_X$ 是\Term{fully-faithful}{全忠实}的。

\begin{SGAresult}{命题}{5.3}
设 $S,Z$ 如5.1，$P$ 是 $S$ 上的\Term{smooth}{光滑}\Term{group-scheme}{群概形}。考虑 $P$ 由 $Z_S$ 给出的\Term{group-scheme}{群概形}\Term{extension}{扩张} $E$。则从这些\Term{extension}{扩张}的\Term{category}{范畴}到 $P_\eta$ 由 $Z_\eta$ 给出的\Term{group-scheme}{群概形}\Term{extension}{扩张}\Term{category}{范畴}的\Term{functor}{函子} $E\mapsto E_\eta$ 是
\SourcePage{290}
\Term{equivalence-of-categories}{范畴等价}。若 $P,Z$ 都交换，此\Term{functor}{函子}也诱导相应的\Term{commutative-extension}{交换扩张}\Term{category}{范畴}之间的\Term{equivalence-of-categories}{等价}。
\end{SGAresult}
\begin{proof}
首先注意，由5.2，$P$ 由 $Z_S$ 给出的每个群\Term{sheaf}{层}\Term{extension}{扩张}（或 $P_\eta$ 由 $Z_\eta$ 给出的\Term{extension}{扩张}）都产生一个平凡\Term{torsor}{挠子}，从而更是\Term{representable}{可表}的。因此，5.3中的\Term{category}{范畴}也\Term{solvable}{可解}释为群\Term{sheaf}{层}的\Term{extension}{扩张}\Term{category}{范畴}。可应用第VII讲1.1.6及1.2中给出的描述：用 $P,P\times P,P\times P\times P$（相应地，$P_\eta,P_\eta\times P_\eta,P_\eta\times P_\eta\times P_\eta$）上，在由 $(Z,Z)$ 经\Term{base-change}{基变换}所得的群对下的\Term{bitorsor}{双挠子}来描述。在\Term{commutative-extension}{交换扩张}的情形应用5.2，在\Term{generic}{一般}情形应用5.2.1，将它们用于 $S$ 上\Term{smooth}{光滑}的\Term{scheme}{概形} $P,P\times P,P\times P\times P$，就得到所宣称的结论。完全类似地得到以下结果。
\end{proof}
\begin{SGAresult}{推论}{5.4}
设 $S,Z$ 如5.1，$P,Q$ 是 $S$ 上的两个\Term{smooth}{光滑}\Term{group-scheme}{群概形}。假设 $P,Q,Z$ 都交换，考虑 $(P,Q)$ 由 $Z_S$ 给出的\Term{biextension}{双扩张} $E$。则从这些\Term{biextension}{双扩张}的\Term{category}{范畴}到 $(P_\eta,Q_\eta)$ 由 $Z_\eta$ 给出的\Term{biextension}{双扩张}\Term{category}{范畴}的\Term{functor}{函子} $E\mapsto E_\eta$ 是\Term{equivalence-of-categories}{范畴等价}。
\end{SGAresult}
上述结果使我们归结到基\Term{scheme}{概形}为域的\Term{spectrum}{谱}时，确定由\Term{torsion-free}{无挠}\Term{constant-group}{常值群}给出的\Term{extension}{扩张}或\Term{biextension}{双扩张}\Term{category}{范畴}的问题。此时有如下结论。

\SourcePage{291}
\begin{SGAresult}{命题}{5.5}
设 $k$ 是域，$Z$ 是\Term{torsion-free}{无挠}群，$P$ 是\Term{locally-finite-type}{局部有限型}的 $k$ \Term{group-scheme}{群概形}，$P^0$ 是\Term{neutral-component}{单位分支}，$\Phi=P/P^0$ 是\Term{connected-component-group}{连通分支群}（SGA 3，第VI$_A$讲4）。则：

（i）从 $\Phi$ 由 $Z_k$ 给出的\Term{extension}{扩张}（相应地，\Term{commutative-extension}{交换扩张}）\Term{group-scheme}{群概形}\Term{category}{范畴}到 $P$ 由 $Z_k$ 给出的\Term{extension}{扩张}（相应地，\Term{commutative-extension}{交换扩张}）\Term{group-scheme}{群概形}\Term{category}{范畴}的自然\Term{functor}{函子}是\Term{equivalence-of-categories}{范畴等价}。特别地，若 $P$ \Term{connected}{连通}，则 $P$ 由 $Z_k$ 给出的每个\Term{extension}{扩张}都平凡，且这些\Term{extension}{扩张}的\Term{category}{范畴}是\Term{rigid-category}{刚性的}。

（ii）假设 $\Phi$ 是\Term{torsion-group}{挠群}（例如 $P$ 为\Term{finite-type}{有限型}），且 $P,Z$ 交换。则 $P$ 由 $Z_k$ 给出的\Term{extension}{扩张}\Term{category}{范畴}是\Term{rigid-category}{刚性的}，即该\Term{category}{范畴}中每个\Term{object}{对象}的任一\Term{automorphism}{自同构}都平凡。再假设存在 $n>0$ 使 $n\cdot\mathrm{id}_\Phi=0$，例如 $\Phi$ \Term{finite}{有限}，比如 $Z$ 为\Term{finite-type}{有限型}，则 $P$ 由 $G_k$ 给出的\Term{commutative-extension}{交换扩张}类群由\Term{canonical-isomorphism}{典范同构}
\[
\operatorname{Ext}^1(P,Z_k)\xleftarrow{\sim}\operatorname{Hom}\bigl(\Phi,(Z\otimes_{\mathbf Z}\mathbf Q)/Z\bigr)\tag{5.5.1}
\]
给出；这个\Term{isomorphism}{同构}由\Term{exact-sequence}{正合列}
\[
0\longrightarrow Z\longrightarrow Z\otimes_{\mathbf Z}\mathbf Q\longrightarrow(Z\otimes_{\mathbf Z}\mathbf Q)/Z\longrightarrow0\tag{5.5.2}
\]
所\Term{derive}{导出}的 $\operatorname{Ext}$ \Term{exact-sequence}{正合列}的\Term{boundary-operator}{上边缘算子}给出。
\end{SGAresult}
\begin{proof}
注意，\Term{canonical}{典范}\Term{morphism}{态射}
\[P\longrightarrow\Phi\]
显然诱导\Term{irreducible-component}{不可约分支}\Term{set}{集合}之间的\Term{isomorphism}{同构}，因为它是\Term{flat}{平坦}且\Term{fiber}{纤维}\Term{irreducible}{不可约}的\Term{morphism}{态射}。
\SourcePage{292}
所以由5.1，从 $Z_\Phi$ 下的\Term{torsor}{挠子}\Term{category}{范畴}到 $Z_P$ 下的\Term{torsor}{挠子}\Term{category}{范畴}的诱导\Term{functor}{函子}是\Term{equivalence-of-categories}{范畴等价}。（回顾，$k$ 上\Term{locally-finite-type}{局部有限型}的\Term{group-scheme}{群概形}是\Term{geometrically-unibranch}{几何单支}的。）同一结论适用于\Term{bitorsor}{双挠子}\Term{category}{范畴}以及\Term{morphism}{态射}
\[P\times P\longrightarrow\Phi\times\Phi,\qquad P\times P\times P\longrightarrow\Phi\times\Phi\times\Phi.\]
第VII讲1.1.6对 $P$ 由 $Z_k$ 给出的\Term{extension}{扩张}和 $\Phi$ 由 $Z_k$ 给出的\Term{extension}{扩张}的描述（\Term{commutative-extension}{交换扩张}的情形见第VII讲1.3），便蕴含（i）的断言。现在只考虑\Term{commutative-extension}{交换扩张}。$P$ 由 $G_k$ 给出的一个\Term{extension}{扩张}的\Term{automorphism-group}{自同构群}为
\[\operatorname{Hom}(P,G_k)\xrightarrow{\sim}\operatorname{Hom}(\Phi,G_k),\]
若 $\Phi$ 是\Term{torsion-group}{挠群}，这个群显然为零。同样，若 $V=G\otimes_{\mathbf Z}\mathbf Q$，且 $\Phi$ \Term{finite}{有限}，则
\[\operatorname{Ext}^i(\Phi,V_k)=0\qquad\text{对每个 }i\in\mathbf Z,\]
因为这些群都是被整数 $n>0$ \Term{annihilate}{零化}的 $\mathbf Q$ \Term{vector-space}{向量空间}。由此立即得到\Term{isomorphism}{同构}（5.5.1）。
\end{proof}
本质上用同样的方法可证明以下结论。
\begin{SGAresult}{推论}{5.6}
设 $k$ 是域，$Z$ 是\Term{torsion-free}{无挠}\Term{commutative-group}{交换群}，$P,Q$ 是 $k$ 上两个\Term{locally-finite-type}{局部有限型}\Term{commutative-group-scheme}{交换群概形}，$\Phi=P/P^0$、$\Psi=Q/Q^0$ 是\Term{neutral-component}{单位分支}群。则：

（i）从 $(\Phi,\Psi)$ 由 $Z_k$ 给出的\Term{biextension}{双扩张}\Term{category}{范畴}
\SourcePage{293}
到 $(P,Q)$ 由 $Z_k$ 给出的\Term{biextension}{双扩张}\Term{category}{范畴}的自然\Term{functor}{函子}是\Term{equivalence-of-categories}{范畴等价}。特别地，若 $P$ 或 $Q$ \Term{connected}{连通}，则 $(P,Q)$ 由 $Z_k$ 给出的\Term{biextension}{双扩张}\Term{category}{范畴}\Term{equivalence-of-categories}{等价}于\Term{point-category}{单点范畴}。

（ii）假设 $\Phi$ 或 $\Psi$ 是\Term{torsion-group}{挠群}。则 $(P,Q)$ 由 $Z_k$ 给出的\Term{biextension}{双扩张}\Term{category}{范畴}是\Term{rigid-category}{刚性的}。再假设存在整数 $n>0$，使 $n\cdot\mathrm{id}_\Phi=0$ 或 $n\cdot\mathrm{id}_\Psi=0$，则 $(P,Q)$ 由 $Z_k$ 给出的\Term{biextension}{双扩张}类群由\Term{canonical-isomorphism}{典范同构}
\[
\operatorname{Biext}^1(P,Q;G_k)\xleftarrow{\sim}\operatorname{Hom}\bigl(\Phi\otimes\Psi,(Z\otimes_{\mathbf Z}\mathbf Q)/Z\bigr)\tag{5.6.1}
\]
给出；此\Term{isomorphism}{同构}由\Term{exact-sequence}{正合列}（5.5.2）所\Term{derive}{导出}的 $\operatorname{Biext}^i$（$i=0,1$）\Term{exact-sequence}{正合列}的\Term{boundary-operator}{上边缘算子}给出。
\end{SGAresult}

\textbf{5.7.} 设 $S$ 是\Term{scheme}{概形}，$T$ 是\Term{closed-subscheme}{闭子概形}，$Z$ 是\Term{set}{集合}。取 $S$ 的 $Z$ 个副本 $S_z$（$z\in Z$）的和，沿 $T$ 的补\Term{open-subset}{开集} $U=S-T$ \Term{gluing}{粘合}所得的\Term{scheme}{概形}，称为 $S$ 上在 $T$ 处\Term{truncation}{截断}、取值于 $Z$ 的\Term{truncated-constant-scheme}{截断常值概形}，记为 $Z_{S,T}$。因此有
\[
Z_{S,T}\times_SU\simeq U,\qquad Z_{S,T}\times_ST\simeq Z_T,\tag{5.7.1}
\]
且 $Z_{S,T}$ 是 $S$ 上的\Term{etale}{平展}\Term{scheme}{概形}。对变动的 $S$ \Term{scheme}{概形} $S'$，有关于 $S'$ \Term{functoriality}{函子性}的\Term{isomorphism}{同构}
\[
\operatorname{Hom}_S(S',Z_{S,T})\xrightarrow{\sim}\operatorname{Hom}_T(S'_T,Z_T),\tag{5.7.2}
\]
其中\Term{arrow}{箭头}由“\Term{restriction}{限制}到 $S'_T$”定义，这里利用了（5.7.1）的第二个\Term{isomorphism}{同构}；此\Term{arrow}{箭头}为\Term{bijection}{双射}的验证留给读者。将\Term{inclusion}{包含映射}记为
\SourcePage{294}
\[i:T\longrightarrow S,\]
则\Term{functoriality}{函子性}\Term{isomorphism}{同构}（5.7.2）也可写成\Term{sheaf}{层}的\Term{isomorphism}{同构}
\[
Z_{S,T}\xrightarrow{\sim}i_*(Z_T).\tag{5.7.3}
\]
从这个公式或直接从定义可见，$Z_{S,T}$ 对变动的 $Z$ 具有\Term{functoriality}{函子性}，且此\Term{functor}{函子}与\Term{finite}{有限}\Term{projective-limit}{投射极限}交换。特别地，它把群变成群，把\Term{commutative-group}{交换群}变成\Term{commutative-group}{交换群}；从 $Z_{S,T}$ 所表示的\Term{functor}{函子}的\Term{form}{形式}（2.7.2）也可清楚看出这一点。

还应指出，从定义或（5.7.3）显然可见，$Z_{S,T}$ 的形成与任意\Term{base-change}{基变换} $f:S'\to S$ 交换，后者给出关于 $Z$ 的\Term{functor}{函子}之间的\Term{isomorphism}{同构}
\[
Z_{S,T}\times_SS'\xrightarrow{\sim}Z_{S',T'},\qquad T'=T\times_SS'=T_{S'}.\tag{5.7.4}
\]
\textbf{5.8.} 现在研究形如 $Z_{S,T}$ 的群下的\Term{torsor}{挠子} $E$，其中 $Z$ 是一个群，为简便起见这里假设它交换。从 $Z_{S,T}$ 下的\Term{torsor}{挠子}\Term{category}{范畴}到 $Z_T$ 下的\Term{torsor}{挠子}\Term{category}{范畴}有自然\Term{functor}{函子}
\[
E\longmapsto E_T.\tag{5.8.0}
\]
我断言，这个\Term{functor}{函子}是\Term{equivalence-of-categories}{范畴等价}。事实上，它\Term{fully-faithful}{全忠实}，即对 $G_{S,T}$ 下的两个\Term{torsor}{挠子} $E,F$，映射 $\operatorname{Hom}(E,F)\to\operatorname{Hom}(E_T,F_T)$ 为\Term{bijection}{双射}；由 fpqc \Term{descent}{下降}，这立即化为 $E,F$
\SourcePage{295}
都平凡的情形，从而化为 $E=F=G_{S,T}$，此时断言正是（5.7.2）对 $S'=S$ 的\Term{bijection}{双射}性。再证明（5.8.0）\Term{essentially-surjective}{本质满射}。设 $E_T$ 是 $Z_T$ 下的\Term{torsor}{挠子}。在5.1的证明中已指出，$E_T$ \Term{representable}{可表}，且对\Term{etale-topology}{平展拓扑}\Term{locally-trivial}{局部平凡}。因此只须证明
\[\mathrm H^1(S_{\text{ét}},Z_{S,T})\longrightarrow\mathrm H^1(T_{\text{ét}},Z_T)\]
为\Term{bijection}{双射}。这来自（5.7.3）及等式
\[\mathrm R^1i_{\text{ét}*}(Z_T)=0,\]
后者是 SGA 4，第VIII讲5.6的特例。

\textbf{5.8.1.} 现在设 $P$ 是 $S$ 上的\Term{commutative-group-scheme}{交换群概形}，我们要确定 $P$ 由 $Z_{S,T}$ 给出的\Term{commutative-group}{交换群}\Term{sheaf}{层}\Term{extension}{扩张}\Term{category}{范畴}。为此，使用第VII讲1.1.6、1.3中的描述：以 $P,P\times P,P\times P\times P$ 上，在 $Z_{S,T}$ 的\Term{inverse-image}{逆像}群下的\Term{torsor}{挠子}来描述。由（5.7.4），这些\Term{inverse-image}{逆像}仍与 $Z_{S,T}$ 同形，因此可将前一结果用于这些\Term{torsor}{挠子}的描述；它表明，它们对应于
\[P_T,\qquad P_T\times_TP_T,\qquad P_T\times_TP_T\times_TP_T\]
上在 $Z_T$ 的\Term{inverse-image}{逆像}群下的\Term{torsor}{挠子}。由此得到以下结果。
\begin{SGAresult}{命题}{5.9}
设 $S$ 是\Term{scheme}{概形}，$T$ 是\Term{closed-subscheme}{闭子概形}，$G$ 是\Term{commutative-group}{交换群}，$G_{S,T}$ 是5.7中定义的 $S$ 上相应的\Term{truncated-constant-group-scheme}{截断常值群概形}，$P$ 是 $S$ 上的\Term{commutative-group-scheme}{交换群概形}。则\Term{restriction}{限制}到 $T$ 的\Term{functor}{函子}诱导
\SourcePage{296}
$P$ 由 $Z_{S,T}$ 给出的\Term{commutative-group}{交换群}\Term{sheaf}{层}\Term{extension}{扩张}\Term{category}{范畴}与 $P_T$ 由 $Z_T$ 给出的\Term{commutative-group}{交换群}\Term{sheaf}{层}\Term{extension}{扩张}\Term{category}{范畴}之间的\Term{equivalence-of-categories}{等价}。
\end{SGAresult}
用本质上相同的证明得到以下结果。
\begin{SGAresult}{推论}{5.10}
沿用5.9的记号，再设 $Q$ 是 $S$ 上的\Term{commutative-group-scheme}{交换群概形}。则\Term{restriction}{限制}到 $T$ 的\Term{functor}{函子}诱导 $(P,Q)$ 由 $Z_{S,T}$ 给出的\Term{biextension}{双扩张}\Term{category}{范畴}与 $(P_T,Q_T)$ 由 $Z_T$ 给出的\Term{biextension}{双扩张}\Term{category}{范畴}之间的\Term{equivalence-of-categories}{等价}。
\end{SGAresult}
